% !TEX TS-program =lualateX
\DocumentMetadata{}
\documentclass[12pt]{article}
\usepackage[headheight=35pt,tmargin=2.5cm,bmargin=3.5cm,lmargin=3cm,rmargin=3cm]{geometry}
\usepackage[colorlinks,linkcolor=blue,citecolor=blue]{hyperref}
\usepackage[spanish,es-noquoting,english]{babel}
\usepackage{graphicx}
%\usepackage[utf8]{inputenc} 
\usepackage{amsmath}
\usepackage{amssymb}
\usepackage{stmaryrd}
\usepackage{textcomp}
\usepackage{fontspec}      % support opentype text fonts
\usepackage[math-style=upright,mathrm=sym,mathbf=sym]{unicode-math}
                           % support opentype math fonts, including bold
%\usepackage[spanish,es-noquoting]{babel}
%\usepackage[T1]{fontenc}
%\usepackage[expert,uprightgreek,altbullet]{lucidabr}
\usepackage{pstricks}
\usepackage{pst-node}
\usepackage{pst-jtree}
\usepackage{pstricks-add}
\usepackage{pst-labo}
\usepackage{pst-coil}
\usepackage[normalem]{ulem}
\usepackage{expex}
\usepackage{phonrule}
\usepackage[most]{tcolorbox}
\usepackage{authblk}
\usepackage{fancyhdr}
%\usepackage{tipa} 
\usepackage[autostyle]{csquotes}
\usepackage{ctable,booktabs}
\usepackage{multirow}
\usepackage{makecell}
\usepackage{enumitem}
\usepackage{parskip}
\usepackage{array}
\usepackage{xcolor}
\usepackage{microtype}

\usepackage[backend=biber, style=authoryear]{biblatex}
\addbibresource{cajal-xmas-2026.bib}

\gathertags

%Open type fonts
\defaultfontfeatures{Scale=.92}
\setmainfont[% main rm
  Ligatures=TeX,
      ItalicFont=LucidaBrightOT-Italic.otf,
        BoldFont=LucidaBrightOT-Demi.otf,
  BoldItalicFont=LucidaBrightOT-DemiItalic.otf,
  ]{LucidaBrightOT.otf}
  %
  \setsansfont[% main sans
  Ligatures=TeX,
      ItalicFont=LucidaSansOT-Italic.otf,
        BoldFont=LucidaSansOT-Demi.otf,
  BoldItalicFont=LucidaSansOT-DemiItalic.otf,
  ]{LucidaSansOT.otf}
%
\setmonofont[% main typewriter
      ItalicFont=LucidaSansTypewriterOT-Oblique.otf,
        BoldFont=LucidaSansTypewriterOT-Bold.otf,
  BoldItalicFont=LucidaSansTypewriterOT-BoldOblique.otf,
  ]{LucidaSansTypewriterOT.otf}
%
\setmathfont{LucidaBrightMathOT.otf}[BoldFont=LucidaBrightMathOT-Demi]
%

\DeclareOldFontCommand{\rm}{\normalfont\rmfamily}{\mathrm}
\everymath={\rm}
\def\TO{\quad$\rightarrow$\quad}

\graphicspath{ {./img/} }
\AtBeginDocument{\lingset{everygla=,aboveglcskip=0pt,aboveglftskip=0pt,glstyle=wrap}}

\renewcommand{\headrulewidth}{0.4pt}
\renewcommand{\footrulewidth}{0.4pt}

\definecolor{accent}{RGB}{140,52,38}
\definecolor{muted}{RGB}{96,88,80}

\newcommand{\es}[1]{\foreignlanguage{spanish}{\textit{#1}}}
\newcommand{\rel}[1]{\textbf{#1}}
\newcommand{\smallcapshead}[1]{{\color{accent}\small\texttt{\MakeUppercase{#1}}}\par\vspace{2pt}{\color{muted}\hrule}\vspace{8pt}}

\title{Closing the center-embedding gap in Spanish:\\can grammatical features do the work of memory?}
\author[1]{Norberto Moreno-Quibén}
\author[2]{Silvia Gumiel-Molina}
\author[3]{Isabel Pérez-Jiménez}
\author[4]{Pilar Pérez-Ocón}
\affil[1,2,3,4]{Departamento de Filología, Comunicación y Documentación, Universidad de Alcalá (UAH)}
\date{{\small {\tt Poster Abstract $\cdot$ Cajal Xmas Meeting 2026 $\cdot$ Instituto Cajal, CSIC, Madrid}}}

\gathertags

\begin{document}
\maketitle

\section*{Abstract}

Center embeddings (clauses nested inside one another in the left branches of the structure) push comprehension to its limits, making them a natural probe for the memory architecture that supports language \parencite{Karlsson2007}. Their signature is the missing-VP illusion: when one of three required verbs is deleted (see below for examples), readers accept the ungrammatical sentence more readily than the grammatical one, as if a stored prediction had silently decayed \parencite{Lewis2005,Hahn2022}. The illusion was recently demonstrated in Spanish \parencite{Paneda2024}. However, illusion magnitude tracked the statistical rarity of the verb sequence, while individual working-memory span predicted nothing. Storage capacity, on its own, appears to be the wrong variable.

Spanish exhibits relativizers that differ in morphological content while occupying identical syntactic positions: featureless \es{que} versus \es{el cual} and \es{quien}, which carry gender, number and animacy. If the illusion arises because degraded memory traces are reconstructed ambiguously, then feature-bearing relativizers should supply distinctive retrieval cues and reduce the illusion, with no change in sentence length, structure, or embedding depth. Featureless \es{que} should not.

We will test this in 120 speakers, combining speeded acceptability judgements, eye-tracking and EEG, and predict a graded reduction in illusion rate with feature richness, an attenuated sustained anterior negativity during maintenance, and a restored P600 at the missing verb.

The result would place grammatical features, not capacity, at the centre of how the brain holds structure open---and hands a broadly neurobiological question, what makes a memory cue distinctive, a linguistically precise manipulation.

\vspace{6pt}
\textbf{Keywords:} sentence processing; cue-based retrieval; working memory; morphosyntactic features; ERP; Spanish

\vspace{14pt}
\smallcapshead{Supporting materials}

\textit{The critical manipulation, on one item. All conditions are double center embeddings of identical depth; the featural content of the relativizer is crossed with grammaticality (VP2 \es{entrevistó} present vs.\ deleted).}

\pex<stim-que>
\a<que-gram> \foreignlanguage{spanish}{La política \rel{a la que} el periodista \rel{al que} los reporteros interrumpieron entrevistó dimitió ayer.} {\color{muted}\small\textsc{[Grammatical]}}
\a<que-ungram> \ljudge{*} \foreignlanguage{spanish}{La política \rel{a la que} el periodista \rel{al que} los reporteros interrumpieron \sout{entrevistó} dimitió ayer.} {\color{muted}\small\textsc{[Ungrammatical, missing VP2]}}\\
{\color{muted}\small\textsc{Baseline} --- \es{que}, gender/number only via the article; the condition used in all prior work.}
\xe

\pex<stim-cual>
\a<cual-gram> \foreignlanguage{spanish}{La política \rel{a la cual} el periodista \rel{al cual} los reporteros interrumpieron entrevistó dimitió ayer.} {\color{muted}\small\textsc{[Grammatical]}}
\a<cual-ungram> \ljudge{*} \foreignlanguage{spanish}{La política \rel{a la cual} el periodista \rel{al cual} los reporteros interrumpieron \sout{entrevistó} dimitió ayer.} {\color{muted}\small\textsc{[Ungrammatical, missing VP2]}}\\
{\color{muted}\small\textsc{Feature-bearing} --- \es{el cual} inflects for gender and number on the pronoun itself.}
\xe

\pex<stim-quien>
\a<quien-gram> \foreignlanguage{spanish}{La política \rel{a quien} el periodista \rel{a quien} los reporteros interrumpieron entrevistó dimitió ayer.} {\color{muted}\small\textsc{[Grammatical]}}
\a<quien-ungram> \ljudge{*} \foreignlanguage{spanish}{La política \rel{a quien} el periodista \rel{a quien} los reporteros interrumpieron \sout{entrevistó} dimitió ayer.} {\color{muted}\small\textsc{[Ungrammatical, missing VP2]}}\\
{\color{muted}\small\textsc{Animacy-bearing} --- \es{quien} restricts the antecedent to a human referent.}
\xe

{\small Differential object marking on \es{a} blocks a subject-relative reading throughout, so no condition is ambiguous at the relativizer.\par}

\vspace{14pt}
\smallcapshead{References}

\nocite{*}
\printbibliography[heading=none]

\end{document}
